\answer{ex:15:1}{(a)~$2\cdot10^5$ por neurônio, $6\cdot10^{12}$ no circuito. (b)~$\approx8$~bit/s por fio, pelo menos $2.5\cdot10^4$~s $\approx7$~h.}
\answer{ex:15:2}{Passo $2$: correlação entre vizinhos $0.943$, $\lambda$ de $4.18$ a $7.3\cdot10^{-4}$, razão $\approx5.7\cdot10^3$; passo $4$: correlação $0.8$, razão $\approx94$.}
\answer{ex:15:3}{(a)~$\lambda=0.988$ e $0.808$: $82$ e $4.7$ movimentos. (b)~$x_s^*=0.690$, $x_f^*=0.172$, soma $0.862$.}
\answer{ex:15:4}{(a)~$116$ movimentos. (b)~$19$. Um modelo com um só processo, com soma zero, não dá economia.}
\answer{ex:15:5}{(a)~$G=50$~s$^{-1}$, contra o limite de $10.5$~s$^{-1}$ sem modelo. (b)~Não: com $G=50$~s$^{-1}$, o atraso crítico é $d_c\approx103\ms$; com $d=150\ms$ é preciso $\hat k>0.445k$ (para quaisquer $G$ e $d$, $\hat k>k/2$).}
\answer{ex:15:6}{O erro cai abaixo de $0.5\%$ em $6$--$30$~s, com piso de $\approx(6$--$9)\cdot10^{-4}$, sendo $7.5\cdot10^{-4}$ com os melhores pesos; com $\eta=50$, os pesos ainda diferem dos melhores em até $0.2$ ao longo das direções mal condicionadas de $C$ (exercício~\ref{ex:15:2}).}
\answer{ex:15:7}{$\mathbf B=A\mathbf K$; $q_f/R\approx0.035$, $q_s/R\approx8.6\cdot10^{-4}$; com $4q_s$, $B_f\approx0.088$, $B_s\approx0.047$: o processo lento aprende mais que o dobro mais depressa, e o rápido, mais devagar.}
